% 第 9 章 Operating Conditions
\bichapter{工作条件}{Operating Conditions}
\label{chap:opcond}

半导体器件参数随工艺、电压与温度（PVT）条件变化。ASIC 厂商通常在多种条件下表征器件参数，并在逻辑库中以不同工作条件集给出参数值。时序分析所用工作条件会影响分析结果。关于为时序分析指定工作条件，见：
\begin{itemize}
  \item 工作条件（Operating Conditions）
  \item 工作条件分析模式（Operating Condition Analysis Modes）
  \item 最小与最大延时计算（Minimum and Maximum Delay Calculations）
  \item 指定分析模式（Specifying the Analysis Mode）
  \item 用两个库进行分析（Using Two Libraries for Analysis）
  \item 延时降额（Derating Timing Delays）
  \item 时钟再汇聚悲观移除（Clock Reconvergence Pessimism Removal）
  \item 时钟片上变异悲观缩减（Clock On-Chip Variation Pessimism Reduction）
\end{itemize}

% ============================================================
\section{工作条件}
\subsection*{Operating Conditions}

集成电路在不同工作条件下呈现不同性能特征：制造工艺变化、电源电压与温度。逻辑库定义这些参数的标称值，并在该条件下给出延时信息。

一组工作条件包含下列值：

\begin{table}[htbp]
\centering
\caption{工作条件参数}
\begin{tabular}{@{}p{3.2cm}p{9.5cm}@{}}
\toprule
工作条件 & 说明 \\
\midrule
工艺降额因子（Process derating factor） & 与制造工艺变化导致的器件参数缩放相关。工艺数小于标称值通常导致更小延时。 \\
环境温度（Ambient temperature） & 芯片温度影响器件延时，取决于环境温度、功耗、封装类型与冷却方式等。 \\
电源电压（Supply voltage） & 较高电源电压通常导致更小延时。 \\
互连模型类型（Interconnect model type） & 定义 PrimeTime 在布局前分析中估计 net 电容与电阻所用的 RC 树拓扑。 \\
\bottomrule
\end{tabular}
\end{table}

每条时序弧的延时信息在标称工艺、温度与电压下指定。若实际工作条件不同，PrimeTime 应用缩放因子以反映这些变化。许多库对工艺、温度与电压使用线性缩放。

若逻辑库包含缩放单元信息，可包含特定工作条件的精确延时表或系数。对非线性缩放库单元，此方法可非常准确。更多信息见 Library Compiler 与 Design Compiler 文档。

可用单组工作条件进行分析（setup 与 hold），或指定最小与最大条件。若未在设计上设置工作条件，PrimeTime 使用主库中的默认工作条件集（若存在），或主库的标称值。

\subsection{互连模型类型}
\subsubsection*{Interconnect Model Types}

PrimeTime 在计算布局前设计的 net 延时时使用互连模型信息（当无反标 net 延时与寄生信息时）。总电阻与电容相同但 RC 树拓扑不同的两条 net，pin 到 pin 延时可能不同。

本主题提供互连模型类型的背景信息。无法在 PrimeTime 中修改这些类型。更多信息见 Library Compiler 文档。

互连模型由每个逻辑库工作条件集中的 \texttt{tree\_type} 规范定义。\texttt{tree\_type} 指示连线电阻与电容拓扑类型：\texttt{best\_case\_tree}、\texttt{worst\_case\_tree} 或 \texttt{balanced\_tree}。例如：
\begin{lstlisting}
operating_conditions(BEST) {
      process     : 1.1;
      temperature : 11.0;
      voltage     : 4.6;
      tree_type   : "best_case_tree";
 }
operating_conditions(TYPICAL) {
      process     : 1.3;
      temperature : 31.0;
      voltage     : 4.6;
      tree_type   : "balanced_tree";
}
operating_conditions(WORST) {
     process      : 1.7;
     temperature : 55.0;
     voltage      : 4.2;
     tree_type    : "worst_case_tree";
}
\end{lstlisting}

若逻辑库未定义树类型，PrimeTime 使用 \texttt{balanced\_tree} 模型。下图示出树类型模型网络。

\figplaceholder{Figure 89: RC interconnect topologies for fanout of N}{扇出为 N 的 RC 互连拓扑}{fig:oc-tree}

\subsection{设置工作条件}
\subsubsection*{Setting Operating Conditions}

要为时序分析指定工艺、温度与电压条件，使用 \cmd{set\_operating\_conditions} 命令。

所指定工作条件必须在指定库或 link path 中的库内定义。要为库创建自定义工作条件，使用 \cmd{create\_operating\_conditions}。使用 \cmd{report\_lib} 获取逻辑库中可用工作条件列表后再使用 \cmd{set\_operating\_conditions}。

将 tech\_lib 库中的 WCCOM 设为单一工作条件：
\begin{lstlisting}
pt_shell> set_operating_conditions WCCOM -library tech_lib
\end{lstlisting}

将 WCCOM 设为最大条件、BCCOM 设为最小条件以进行片上变异分析：
\begin{lstlisting}
pt_shell> set_operating_conditions -analysis_type on_chip_variation \
            -min BCCOM -max WCCOM
\end{lstlisting}

未指定库时，PrimeTime 搜索 link path 中所有库。设置工作条件后，可报告或移除工作条件。

\subsection{创建工作条件}
\subsubsection*{Creating Operating Conditions}

逻辑库包含固定的工作条件集。要在库中创建新工作条件，使用 \cmd{create\_operating\_conditions}。这些自定义工作条件可用于当前会话中的设计分析，但不能写入库 \texttt{.db} 文件。

查看库中定义的工作条件用 \cmd{report\_lib}；在当前设计上设置工作条件用 \cmd{set\_operating\_conditions}。

在 tech\_lib 库中创建名为 WC\_CUSTOM 的新工作条件：
\begin{lstlisting}
pt_shell> create_operating_conditions -name WC_CUSTOM \
            -library tech_lib -process 1.2 \
            -temperature 30.0 -voltage 2.8 \
            -tree_type worst_case_tree
\end{lstlisting}

\subsection{工作条件信息命令}
\subsubsection*{Operating Condition Information}

下列命令用于报告、移除或重置工作条件信息：

\begin{table}[htbp]
\centering
\caption{工作条件相关命令}
\begin{tabular}{@{}ll@{}}
\toprule
命令 & 操作 \\
\midrule
\cmd{report\_design} & 列出设计的工作条件设置 \\
\cmd{remove\_operating\_conditions} & 从当前设计移除工作条件 \\
\cmd{reset\_design} & 将工作条件重置为默认并移除所有用户指定数据（如时钟、输入/输出延时） \\
\bottomrule
\end{tabular}
\end{table}

% ============================================================
\section{工作条件分析模式}
\subsection*{Operating Condition Analysis Modes}

半导体器件参数可随制造工艺、工作温度与电源电压等条件变化。PrimeTime 中 \cmd{set\_operating\_conditions} 指定分析工作条件，使 PrimeTime 使用逻辑库中相应的参数值集。

PrimeTime 提供下列设置时序分析工作条件的方法：
\begin{itemize}
  \item \textbf{单一工作条件模式}：基于一组工艺、温度与电压条件，对整个设计使用单组延时参数。
  \item \textbf{片上变异（OCV）模式}：执行保守分析，允许最小与最大延时同时应用于不同路径。Setup 检查对 launch 时钟路径与数据路径用最大延时，对 capture 时钟路径用最小延时。Hold 检查相反。
  \item \textbf{高级片上变异（AOCV）模式}：根据路径逻辑深度与物理距离等指标确定降额因子。
  \item \textbf{参数化片上变异（POCV）模式}：将延时、要求时间与 slack 计算为统计分布。
\end{itemize}

下表示出各工作条件分析模式下 setup 与 hold 检查所用的时钟到达时间、延时、工作条件与延时降额。

\begin{table}[htbp]
\centering
\caption{表 16：Setup 与 Hold 检查所用的时序参数（续表见下文）}
\small
\begin{tabular}{@{}p{1.8cm}p{1.2cm}p{3.2cm}p{3.2cm}p{3.2cm}@{}}
\toprule
分析模式 & 时序检查 & Launch 时钟路径 & 数据路径 & Capture 时钟路径 \\
\midrule
单一工作条件 & Setup & 晚时钟，时钟路径最大延时，单一工作条件（无降额） & 最大延时，单一工作条件（无降额） & 早时钟，时钟路径最小延时，单一工作条件（无降额） \\
单一工作条件 & Hold & 早时钟，时钟路径最小延时，单一工作条件（无降额） & 最小延时，单一工作条件（无降额） & 晚时钟，时钟路径最大延时，单一工作条件（无降额） \\
\bottomrule
\end{tabular}
\end{table}

\begin{table}[htbp]
\centering
\caption{表 16（续）：OCV 模式}
\small
\begin{tabular}{@{}p{1.8cm}p{1.2cm}p{3.2cm}p{3.2cm}p{3.2cm}@{}}
\toprule
分析模式 & 时序检查 & Launch 时钟路径 & 数据路径 & Capture 时钟路径 \\
\midrule
OCV & Setup & 晚时钟，时钟路径最大延时，晚降额，最坏工作条件 & 最大延时，晚降额，最坏工作条件 & 早时钟，时钟路径最小延时，早降额，最好工作条件 \\
OCV & Hold & 早时钟，时钟路径最小延时，早降额，最好工作条件 & 最小延时，早降额，最好工作条件 & 晚时钟，时钟路径最大延时，晚降额，最坏工作条件 \\
\bottomrule
\end{tabular}
\end{table}

% ============================================================
\section{最小与最大延时计算}
\subsection*{Minimum and Maximum Delay Calculations}

\cmd{set\_operating\_conditions} 定义时序分析的工作条件并指定分析类型（单一或片上变异）。工作条件必须在指定库或库对中定义。

默认情况下，PrimeTime 在单一工作条件集下分析（单一工作条件模式）。此模式下需多次分析运行以处理多种工作条件。通常至少需分析两种工作条件以确保无时序违例：最好情况（最小路径报告）用于 hold 检查，最坏情况（最大路径报告）用于 setup 检查。

在片上变异（OCV）模式下，PrimeTime 执行保守分析，允许最小与最大延时同时应用于不同路径。Setup 检查对 launch 时钟路径与数据路径用最大延时，对 capture 时钟路径用最小延时；Hold 检查相反。

OCV 模式下，当路径段被 launch 与 capture 时钟路径共享时，该段可能同时被视为具有两种不同延时。未考虑共享路径段会导致悲观分析。可通过修正获得更准确分析，见「时钟再汇聚悲观移除」。

最小-最大分析考虑下列设计参数的最小与最大值：
\begin{itemize}
  \item 输入与输出外部延时
  \item 自 SDF 反标的延时
  \item 端口连线负载模型
  \item 端口扇出数
  \item Net 电容与电阻
  \item Net 连线负载模型
  \item 时钟 latency 与转换时间
  \item 输入端口驱动单元
\end{itemize}

例如计算最大延时时，PrimeTime 使用最长路径、最坏工作条件、最晚到达时钟边沿、最大单元延时、最长转换时间等。

可用单库（指定最小与最大工作条件）或双库（一库最好、一库最坏）进行最小-最大分析，见「用两个库进行分析」。

启用最小-最大分析的方法：
\begin{itemize}
  \item 用 \cmd{set\_operating\_conditions} 设置最小与最大工作条件：
\begin{lstlisting}
pt_shell> set_operating_conditions -analysis_type on_chip_variation \
          -min BCCOM -max WCCOM
\end{lstlisting}
  \item 读入 SDF 并读取最小与最大延时值：
\begin{lstlisting}
pt_shell> read_sdf -analysis_type on_chip_variation my_design.sdf
\end{lstlisting}
\end{itemize}

\subsection{最小-最大单元与 Net 延时值}
\subsubsection*{Minimum-Maximum Cell and Net Delay Values}

每条时序弧可有最小与最大延时以反映工作条件变化。指定方式：
\begin{itemize}
  \item 从一个或两个 SDF 文件反标延时
  \item 由 PrimeTime 计算延时
\end{itemize}

从 SDF 反标：SDF 三元组的最小/最大、两个 SDF 文件、或上述方式加最大/最小乘数。

由 PrimeTime 计算：单一工作条件加时序降额、两种工作条件（最好/最坏）建模 OCV、上述加最小/最大降额、单元与 net 分别降额、不同时序检查分别降额。

\begin{table}[htbp]
\centering
\caption{表 17：来自 SDF 三元组数据的最小-最大延时}
\small
\begin{tabular}{@{}p{2.8cm}p{4.5cm}p{2cm}p{2cm}@{}}
\toprule
分析模式 & 基于工作条件的延时 & 一个 SDF & 两个 SDF \\
\midrule
单一工作条件 Setup & 工作条件下最大数据；工作条件下最小 capture 时钟 & \texttt{a:b:c} & --- \\
单一工作条件 Hold & 工作条件下最小数据；工作条件下最大 capture 时钟 & \texttt{a:b:c} & --- \\
OCV Setup & 最坏最大数据；最好最小 capture 时钟 & \texttt{a:b:c} & SDF1/SDF2 \\
OCV Hold & 最好最小数据；最坏最大 capture 时钟 & \texttt{a:b:c} & SDF1/SDF2 \\
\bottomrule
\end{tabular}
\end{table}

PrimeTime 使用粗斜体显示的三元组值。

\subsubsection{Setup 与 Hold 检查}
\paragraph*{Setup and Hold Checks}

单一工作条件与 OCV 下 setup/hold 检查示例见：
\begin{itemize}
  \item 路径延时追踪
  \item 最坏条件下 Setup 时序检查
  \item 最好条件下 Hold 时序检查
\end{itemize}

\paragraph{路径延时追踪}
\subparagraph*{Path Delay Tracing for Setup and Hold Checks}

下图示出 PrimeTime 如何执行 setup 与 hold 检查。

\figplaceholder{Figure 90: Design Example}{设计示例}{fig:oc-path-trace}

\begin{itemize}
  \item Setup 检查（DL2/ck 到 DL2/d）：clock\_path1 最大延时；数据路径（data\_path\_max）最大延时；clock\_path2 最小延时。
  \item Hold 检查：clock\_path1 最小延时；数据路径（data\_path\_min）最小延时；clock\_path2 最大延时。
\end{itemize}

data\_path\_min 与 data\_path\_max 可因组合逻辑中多条拓扑路径而不同。

\paragraph{最坏条件下 Setup 时序检查}
\subparagraph*{Setup Timing Check for Worst-Case Conditions}

\figplaceholder{Figure 91: Setup Check Using Worst-Case Conditions}{用最坏条件的 Setup 检查}{fig:oc-setup-wc}

PrimeTime 按下列方式检查 setup 违例：
\[
\text{clockpath1} + \text{datapathmax} - \text{clockpath2} + \text{setup} \leq \text{clockperiod}
\]
其中 clockpath1 $= 0.8 + 0.6 = 1.4$，datapathmax $= 3.8$，clockpath2 $= 0.8 + 0.65 = 1.45$，setup $= 0.2$。时钟周期至少为 $1.4 + 3.8 - 1.45 + 0.2 = 3.95$。

\paragraph{最好条件下 Hold 时序检查}
\subparagraph*{Hold Timing Check for Best-Case Conditions}

\figplaceholder{Figure 92: Hold Check Using Best-Case Conditions}{用最好条件的 Hold 检查}{fig:oc-hold-bc}

Hold 违例检查：
\[
\text{clockpath1} + \text{datapathmax} - \text{clockpath2} - \text{hold} \geq 0
\]
无 hold 违例，因 $0.8 + 1.6 - 0.85 - 0.1 = 1.45 > 0$。

\paragraph{存在延时变异时的路径追踪}
\subparagraph*{Path Tracing in the Presence of Delay Variation}

OCV 下每个单元或 net 延时存在不确定性。可指定最坏条件下 OCV 为标称延时的 80\%--100\%。

\figplaceholder{Figure 93: OCV for Worst-Case Conditions}{最坏条件的 OCV}{fig:oc-ocv-wc}

此模式下，给定路径的最大延时按最坏情况 100\% 计算，最小按 80\%。Setup 检查：
\[
\text{clockpath1} + \text{datapathmax} - \text{clockpath2} - \text{setup} \leq 0
\]
时钟周期至少为 $1.40 + 3.80 - 1.16 + 0.2 = 4.24$。

\begin{noteBox}
片上变异影响时钟 latency；因此仅在使用传播时钟 latency 时需要。若指定理想时钟 latency，可通过 \cmd{set\_clock\_uncertainty} 增大时钟不确定性以考虑 OCV。
\end{noteBox}

% ============================================================
\section{指定分析模式}
\subsection*{Specifying the Analysis Mode}

不同工作条件下时序分析示例见：
\begin{itemize}
  \item 单一工作条件分析
  \item 片上变异分析
\end{itemize}

\subsection{单一工作条件分析}
\subsubsection*{Single Operating Condition Analysis}

最好条件分析：
\begin{lstlisting}
pt_shell> set_operating_conditions BEST
pt_shell> report_timing -delay_type min
\end{lstlisting}

\figplaceholder{Figure 94: Timing path for one operating condition -- best case}{单一工作条件时序路径 -- 最好情况}{fig:oc-single-best}

最坏条件分析：
\begin{lstlisting}
pt_shell> set_operating_conditions WORST
pt_shell> report_timing -delay_type max
\end{lstlisting}

\figplaceholder{Figure 95: Timing path for one operating condition -- worst case}{单一工作条件时序路径 -- 最坏情况}{fig:oc-single-worst}

\subsection{片上变异分析}
\subsubsection*{On-Chip Variation Analysis}

执行 OCV 分析：
\begin{enumerate}
  \item 指定工作条件：
\begin{lstlisting}
pt_shell> set_operating_conditions \
            -analysis_type on_chip_variation \
            -min MIN -max MAX
\end{lstlisting}
  \item 报告 setup 与 hold：
\begin{lstlisting}
pt_shell> report_timing -delay_type min
pt_shell> report_timing -delay_type max
\end{lstlisting}
\end{enumerate}

\figplaceholder{Figure 96: Timing path reported during OCV analysis}{OCV 分析报告的时序路径}{fig:oc-ocv-path}

\figplaceholder{Figure 97: Early and late arrival time}{早到与晚到时间}{fig:oc-early-late}

\textbf{示例 1}：WCCOM 最坏商业条件下 OCV 20\% 以下：
\begin{lstlisting}
pt_shell> set_operating_conditions -analysis_type \
            on_chip_variation WCCOM
pt_shell> set_timing_derate -early 0.8
pt_shell> report_timing
\end{lstlisting}

\textbf{示例 2}：WCCOM\_scaled 与 WCCOM 之间 OCV：
\begin{lstlisting}
pt_shell> set_operating_conditions -analysis_type \
            on_chip_variation -min WCCOM_scaled -max WCCOM
pt_shell> report_timing
\end{lstlisting}

\textbf{示例 3}：单元延时 OCV 在 SDF 值上下 5\%/10\%；net 延时 2\%/4\%；单元时序检查 setup 上 10\%、hold 下 20\%：
\begin{lstlisting}
pt_shell> set_timing_derate -cell_delay -early 0.90
pt_shell> set_timing_derate -cell_delay -late 1.05
pt_shell> set_timing_derate -net_delay -early 0.96
pt_shell> set_timing_derate -net_delay -late 1.02
pt_shell> set_timing_derate -cell_check -early 0.80
pt_shell> set_timing_derate -cell_check -late 1.10
\end{lstlisting}

% ============================================================
\section{用两个库进行分析}
\subsection*{Using Two Libraries for Analysis}

\cmd{set\_min\_library} 指示 PrimeTime 同时使用两个逻辑库进行最小与最大延时分析，例如：
\begin{itemize}
  \item 最好与最坏工作条件
  \item 乐观与悲观连线负载模型
  \item 最小与最大时序延时
\end{itemize}

用 \cmd{set\_min\_library} 在两库间建立最小/最大关联：一库用于最大延时分析，另一用于最小延时分析。仅最大库应在 link path 中。

使用 \cmd{set\_min\_library} 时，PrimeTime 先查最大库中的库单元，再在最小库中查找匹配（同名、同 pin、同时序弧）。若存在则用于最小分析，否则使用最大库信息。

% ============================================================
\section{延时降额}
\subsection*{Derating Timing Delays}

可对计算延时降额以建模片上变异效应。降额将计算延时乘以指定因子，改变 \cmd{report\_timing} 等报告的延时与 slack。

\cmd{set\_timing\_derate} 指定早/晚延时调整因子及可选作用范围：
\begin{lstlisting}
pt_shell> set_timing_derate -early 0.9
pt_shell> set_timing_derate -late 1.2
\end{lstlisting}

第一条将所有早（最短路径）单元与 net 延时减少 10\%（如 hold 数据路径）；第二条将晚延时增加 20\%（如 setup 数据路径），使分析更保守。

\cmd{set\_timing\_derate} 会隐式切换到 OCV 模式（若尚未处于该模式），等效于：
\begin{lstlisting}
pt_shell> set_operating_conditions -analysis_type on_chip_variation
\end{lstlisting}

OCV 模式下同时对早、晚应用最坏情况调整。Setup 检查对 launch 时钟与数据路径用晚降额，对 capture 时钟路径用早降额；Hold 检查相反。

必须指定 \opt{-early} 或 \opt{-late}。早、晚降额需两条命令。降额因子为浮点数；更保守分析用早降额 $< 1.0$ 或晚降额 $> 1.0$。

报告已设降额用 \cmd{report\_timing\_derate}。在 \cmd{report\_timing} 中用 \opt{-derate} 报告每级增量延时的降额因子。

取消所有降额用 \cmd{reset\_timing\_derate}。

\subsection{降额选项}
\subsubsection*{Derating Options}

默认 \cmd{set\_timing\_derate} 将指定早/晚降额因子应用于全设计所有早/晚时序路径。可选缩小范围：
\begin{itemize}
  \item \opt{-clock}/\opt{-data}：仅时钟/数据路径
  \item \opt{-net\_delay}/\opt{-cell\_delay}：仅 net/单元延时
  \item \opt{-rise}/\opt{-fall}：仅上升/下降边延时
  \item 对象列表：特定 net、单元实例或库单元
  \item \opt{-net\_delay} 与 \opt{-static}/\opt{-dynamic}：仅静态或动态 net 延时分量
  \item \opt{-cell\_check}：单元时序检查约束（hold/removal 用早降额；setup/recovery 用晚降额）
  \item \opt{-aocvm\_guardband}、\opt{-pocvm\_guardband}、\opt{-pocvm\_coefficient\_scale\_factor}：影响 AOCV/POCV 分析
\end{itemize}

在 net 或叶级单元上设置降额时，命令隐式匹配对象类型。层次单元上必须显式指定 \opt{-net\_delay}、\opt{-cell\_delay} 或 \opt{-cell\_check}。

将 \cmd{timing\_use\_constraint\_derates\_for\_pulse\_checks} 设为 \texttt{true} 可使最小脉宽与最小周期约束参与降额。也可用 \cmd{set\_timing\_derate} 的 \opt{-min\_pulse\_width}、\opt{-min\_period} 选项（此时上述变量须为 \texttt{false}，默认）。

\begin{lstlisting}
set_timing_derate -min_pulse_width rise 1.2
set_timing_derate -min_period fall 1.1
report_timing_derate -min_pulse_width
report_timing_derate -min_period
reset_timing_derate -min_pulse_width
reset_timing_derate -min_period
\end{lstlisting}

此功能与 POCV 分析配合；在 \cmd{set\_timing\_derate} 中与 \opt{-pocvm\_guardband}、\opt{-pocvm\_coefficient\_scale\_factor} 联用。

\subsection{冲突的降额设置}
\subsubsection*{Conflicting Derating Settings}

新 \cmd{set\_timing\_derate} 覆盖同范围先前设置的降额值。范围重叠时，范围更窄的命令优先：
\begin{lstlisting}
set_timing_derate -late -cell_delay 1.4 [get_cells U2*]
set_timing_derate -late -cell_delay 1.3 [get_lib_cells libAZ/ND*]
set_timing_derate -late -cell_delay 1.2
set_timing_derate -late -1.1
\end{lstlisting}

\subsection{负延时降额}
\subsubsection*{Derating Negative Delays}

某些情况下延时可为负（如输入转换慢、输出转换快且延时极短时，输出可在输入到达 50\% 触发电平前到达）。串扰导致 net 延时变化也可能类似。

一般调整公式：
\[
\text{delay\_new} = \text{old\_delay} + ((\text{derating\_factor} - 1.0) \times |\text{old\_delay}|)
\]
正延时时简化为 $\text{delay\_new} = \text{old\_delay} \times \text{derating\_factor}$；负延时时为 $\text{delay\_new} = \text{old\_delay} \times (2.0 - \text{derating\_factor})$。

静态与动态 net 延时分量分别降额后再合并，使负 delta 延时在合并到正静态延时前正确降额。

% ============================================================
\section{时钟再汇聚悲观移除}
\subsection*{Clock Reconvergence Pessimism Removal}

时钟再汇聚悲观是：两条不同时钟路径部分共享物理路径段，共享段被假定对一路径为最小延时、对另一路径为最大延时。launch 与 capture 时钟路径使用不同延时时常出现，最常见于 OCV 分析。自动修正称为时钟再汇聚悲观移除（CRPR）。

CRPR 默认使能。可禁用以节省运行时间与内存（牺牲精度、增加悲观）。CRPR 仅减少悲观，故仅能增加报告 slack。若禁用 CRPR 时设计无违例，使能 CRPR 后仍无违例。

禁用 CRPR：在时序分析前将 \cmd{timing\_remove\_clock\_reconvergence\_pessimism} 设为 \texttt{false}。更改此变量会导致完整时序更新。

\subsection{片上变异示例}
\subsubsection*{On-Chip Variation Example}

\begin{lstlisting}
pt_shell> set_operating_conditions -analysis_type \
          on_chip_variation -min MIN -max MAX
pt_shell> set_timing_derate -net_delay -early 0.80
pt_shell> report_timing -delay_type min
pt_shell> report_timing -delay_type max
\end{lstlisting}

\figplaceholder{Figure 98: Clock reconvergence pessimism example}{时钟再汇聚悲观示例}{fig:crpr-ocv}

各延时有最小与最大值（最好/最坏工作条件）。LD2/CP 的 setup 检查将源锁存器时钟路径（CLK 到 LD1/CP）按 100\% 最坏，目标锁存器路径（CLK 到 LD2/CP）按 80\% 最坏。

虽为有效方法，但因 clock path1 与 path2 共享至 U1 输出的时钟树，分析悲观：共享段称为公共部分，最后单元输出为公共点（common point）。Setup 检查假定 U1 同时有 0.64 与 0.80 两种延时，悲观量为 0.16（公共点最晚与最早到达时间之差），即时钟再汇聚悲观。Hold 检查类似。

\subsection{再汇聚逻辑示例}
\subsubsection*{Reconvergent Logic Example}

\figplaceholder{Figure 99: Reconvergent logic in a clock network}{时钟网络中的再汇聚逻辑}{fig:crpr-reconv}

即使无 OCV，时钟网络中的再汇聚逻辑也可能产生时钟再汇聚：馈入多路选择器的两路时钟路径不能同时有效，但分析可能对同一次 setup/hold 同时考虑较短与较长路径。

\subsection{最小脉宽检查示例}
\subsubsection*{Minimum Pulse Width Checking Example}

\cmd{report\_constraint} 检查时钟网络（\cmd{set\_min\_pulse\_width}）与单元输入（逻辑库）的最小脉宽违例。CRPR 提高此检查精度。

\figplaceholder{Figure 100: Minimum pulse width analysis}{最小脉宽分析}{fig:crpr-mpw}

无 CRPR 时最坏脉宽可能极小并违例，因假设上升与下降边同时最坏延时。实际电路中两边延时至少部分相关。使能 CRPR 时，工具将一定 slack 加回最小脉宽计算：
\[
\text{crp} = \min[(M_r - m_r), (M_f - m_f)]
\]
$M_r/m_r$ 为累积最大/最小上升延时，$M_f/m_f$ 为下降。

\subsection{CRPR 报告}
\subsubsection*{CRPR Reporting}

\cmd{report\_timing}、\cmd{report\_constraint}、\cmd{report\_analysis\_coverage}、\cmd{report\_bottleneck} 详细报告会报告 CRPR 调整。示例 \cmd{report\_timing -delay_type max} 输出中含 \texttt{clock reconvergence pessimism} 行。

\subsection{公共点不同跳变的 CRPR 降额}
\subsubsection*{Derating CRPR for Different Transitions at the Common Point}

\cmd{timing\_clock\_reconvergence\_pessimism} 指定 launch 与 capture 时钟路径跳变不匹配时如何找最晚公共点。\texttt{normal}（默认）找最晚拓扑公共点（跳变可不同）；\texttt{same\_transition} 还要求跳变类型相同。

\figplaceholder{Figure 101: Common Point for Different timing\_clock\_reconvergence\_pessimism Settings}{不同 CRPR 变量设置下的公共点}{fig:crpr-common}

\texttt{normal} 适用于上升/下降到达时间高度相关；不高度相关时 \texttt{same\_transition} 保证保守（可能更悲观）分析。

可用 \texttt{normal} 公共点搜索并对公共点跳变不同时的 CRP 降额：
\begin{lstlisting}
set_app_var timing_crpr_different_transition_derate 0.90
\end{lstlisting}

跳变不同时 CRP 减至原值的 90\% 再加回 slack，比默认 100\% 更保守。

POCV 分析中可同时降额 CRP 变化：
\begin{lstlisting}
set_app_var timing_crpr_different_transition_derate 0.90
set_app_var timing_crpr_different_transition_variation_derate 0.90
\end{lstlisting}

公共点跳变相同时无效果；\texttt{timing\_clock\_reconvergence\_pessimism} 为 \texttt{same\_transition} 时也无效果。第一变量默认 1.0（标称不降额），第二默认 0.0（完全消除 CRP 变化）。

最小脉宽示例脚本：
\begin{lstlisting}
set_operating_conditions -analysis_type on_chip_variation
create_clock -period 8 CLK
set_propagated_clock [all_clocks]
set_clock_latency -source -early -1.3 CLK
set_clock_latency -source -late 1.4 CLK
set_annotated_delay -cell -from CT1/A -to CT1/Z 1
set_annotated_delay -cell -from CT2/A -to CT2/Z -min -rise 0.8
set_annotated_delay -cell -from CT2/A -to CT2/Z -max -rise 1.8
set_annotated_delay -cell -from CT2/A -to CT2/Z -min -fall 0.85
set_annotated_delay -cell -from CT2/A -to CT2/Z -max -fall 2.03
\end{lstlisting}

使能 CRPR 时 \cmd{report\_constraint} 路径报告含 \texttt{clock reconvergence pessimism 3.70} 等。

\subsection{CRPR 合并阈值}
\subsubsection*{CRPR Merging Threshold}

为计算效率，相邻点 CRP 差小于阈值时合并路径点。\cmd{timing\_crpr\_threshold\_ps} 以皮秒指定阈值，默认 5 ps。建议在性能与精度间取时钟网络典型门级延时（门延时+net 延时）的一半；设计阶段可用较大值，签核用较小值。

\subsection{CRPR 与串扰分析}
\subsubsection*{CRPR and Crosstalk Analysis}

用 PrimeTime SI 做串扰分析时，时钟路径公共段上因串扰的延时变化对零周期检查可能悲观（同一时钟边沿驱动 launch 与 capture）。其他路径类型不悲观，因不能假定 launch 与 capture 边沿的串扰变化相同。

因此 CRPR 仅在零周期检查时移除公共段上的串扰诱导延时。适用情形包括：标准 hold 检查、Q-bar 接 D 的寄存器、串扰反馈 hold、多周期路径为零的 hold、涉及透明锁存器的某些 setup 检查等。

\subsection{CRPR 与动态时钟到达}
\subsubsection*{CRPR With Dynamic Clock Arrivals}

动态反标（\cmd{set\_clock\_latency}、\cmd{set\_voltage} 设置的动态时钟 latency 与 rail 电压）可导致动态时钟到达。CRPR 与处理串扰延时相同：仅对零周期路径用动态时钟到达计算 CRP；其他路径仅用静态分量。

示例：早源 latency 2.5（静态 3.0 + 动态 $-0.5$）；晚源 latency 5.5（静态 5.0 + 动态 0.5）。静态 CRP $= 2$（5 减 3）；动态 CRP $= 3$（5.5 减 2.5）。

\subsection{透明锁存器边沿考虑}
\subsubsection*{Transparent Latch Edge Considerations}

路径终点为透明锁存器时，PrimeTime 计算两个 CRP 值：公共节点上升与下降各一。开/关时钟边沿各平移对应悲观量。

\figplaceholder{Figure 102: CRPR for latches}{锁存器的 CRPR}{fig:crpr-latch}

开边沿悲观影响非借用路径 slack 并减少借用路径借用时间；关边沿悲观增加最大借用时间并减少 setup 违例量。用 \cmd{report\_crpr} 报告锁存器 CRP 计算。

\subsection{报告 CRPR 计算}
\subsubsection*{Reporting CRPR Calculations}

\cmd{report\_crpr} 报告两寄存器时钟 pin 或端口间的 CRP 计算，含静态/动态条件及实际用于悲观移除的值。需指定 launch/capture pin、时钟与检查类型（setup/hold）：
\begin{lstlisting}
pt_shell> report_crpr -from [get_pins ffa/CP] \
            -to [get_pins ffd/CP] \
            -from_clock CLK -setup
\end{lstlisting}

\cmd{report\_crpr} 与 \cmd{report\_timing} 报告的 CRP 可能略有不同，因 \cmd{report\_timing} 为效率合并相邻点（\cmd{timing\_crpr\_threshold\_ps}，默认 5 ps）。不一致时 \cmd{report\_crpr} 更准确。

% ============================================================
\section{时钟片上变异悲观缩减}
\subsection*{Clock On-Chip Variation Pessimism Reduction}

\cmd{set\_timing\_derate} 可建模 OCV。launch 与 capture 时钟路径有公共段时，若使能 CRPR，会移除公共段延时降额导致的悲观，但默认仅在路径最终 slack 计算后移除，不在计算串扰到达窗口时移除。若通向 aggressor/victim 的时钟路径共享公共段，到达窗口可能过宽，导致边际重叠的串扰情形在准确 CRPR 下本不存在。

\figplaceholder{Figure 103: Clock reconvergence pessimism in crosstalk arrival windows}{串扰到达窗口中的时钟再汇聚悲观}{fig:ocv-xtalk}

无降额时到达窗口宽 1.0，不重叠；降额后公共点前长时钟路径早到 6.0、晚到 9.0，窗口扩至 3.0 并重叠，产生实际电路中不存在的串扰。

为在基于路径分析中获得最佳精度，可选在到达窗口重叠分析中应用 CRPR，移除公共点早/晚到达差异导致的悲观。设置：
\begin{lstlisting}
set_app_var pba_enable_xtalk_delay_ocv_pessimism_reduction true
\end{lstlisting}
默认 \texttt{false}。为 \texttt{true} 且 CRPR 使能时，PrimeTime SI 在计算 aggressor/victim 到达窗口时应用 CRPR，提高精度，代价是额外运行时间。基于路径分析在使用 \cmd{get\_timing\_paths} 或 \cmd{report\_timing} 的 \opt{-pba\_mode} 时发生。CRPR 在 \cmd{timing\_remove\_clock\_reconvergence\_pessimism} 为 \texttt{true} 时使能。
